\chapter{Formalization contract and provenance}
\label{chap:formalization-contract}

This chapter fixes the trust boundary used throughout the Blueprint.  The
mathematical source of record is Lin--Wang--Xu, \emph{On the Last Kervaire
Invariant Problem}, represented by \texttt{aimpaper/main.tex},
\texttt{aimpaper/112.tex}, and the local PDF.  Earlier literature and machine
calculations may enter only through the explicit records below.  In
particular, a citation or a table row is not silently promoted to a project
axiom.

\section{Source-bearing inputs}

% SOURCE: PROJECT_BOUNDARY.md, External inputs; Acceptance criteria.
\begin{definition}[Source locator]
  \label{def:source-ref}
  \notready
  A source locator is a record consisting of an author or owner, a title or
  dataset name, a stable identifier, a locator inside that source (theorem,
  page, table, row, file and line range, or machine-output key), and a source
  class.  The allowed source classes are \emph{paper}, \emph{literature},
  \emph{machine computation}, and \emph{project governance}.  A locator is
  required to identify the proposition actually supplied, rather than merely
  a nearby document.
\end{definition}

% SOURCE: PROJECT_BOUNDARY.md, External inputs, conceptual Exterresult record.
\begin{definition}[External result]
  \label{def:external-result}
  \uses{def:source-ref}
  \notready
  For a proposition $P$, an external result is a pair $(p,\sigma)$ with
  $p:P$ and a literature source locator $\sigma$ whose cited statement is
  $P$.  It is supplied as an explicit argument to every theorem that uses it;
  the project shall not install it as an axiom or an untracked global
  instance.
\end{definition}

% SOURCE: PROJECT_BOUNDARY.md, External inputs, conceptual Exterevidence record.
\begin{definition}[External evidence]
  \label{def:external-evidence}
  \uses{def:source-ref}
  \notready
  For a proposition $P$, external evidence is a record containing evidence
  $p:P$, a source locator, the computational or transcription method, and
  enough indexing metadata to identify the input and output.  It is used for
  finite Ext data, differentials, table entries, and diagram relations.  Like
  an external result, it is always an explicit theorem parameter.
\end{definition}

% SOURCE: PROJECT_BOUNDARY.md, Examples, remarks, and questions.
\begin{definition}[Open proposition]
  \label{def:open-proposition}
  \notready
  An open proposition is a proposition recorded as a formalization target but
  neither assumed nor asserted.  Its Blueprint node has no incoming edge to a
  proved conclusion unless a later theorem explicitly takes a proof of it as
  a hypothesis.
\end{definition}

% SOURCE: PROJECT_BOUNDARY.md, Confirmed design decisions and Axiom policy.
\begin{definition}[KIP126 trust boundary]
  \label{def:kip126-trust-boundary}
  \uses{def:external-result,def:external-evidence,def:open-proposition,def:kervaire-main-input}
  \notready
  The KIP126 development may use Lean's foundational axioms and declarations
  proved in the pinned Mathlib revision.  Paper-internal arguments are project
  proof obligations.  Results from earlier literature are explicit external
  results, while high-stem computations and every appendix row are explicit
  external evidence.  The conditional permanent-cycle and geometric endpoints
  are universally quantified over the explicit main-input record of
  Definition~\ref{def:kervaire-main-input}; none of its fields is installed as
  a global instance.  No project declaration may use \texttt{axiom},
  \texttt{sorry}, or \texttt{admit}.
\end{definition}

The comments beginning with \texttt{SOURCE} in later chapters give the primary
locator.  Comments beginning with \texttt{VERIFIED-IN} record a second local
copy that was checked, but never replace the primary source.  A
\texttt{mathlibok} node has an exact declaration in Mathlib v4.32.2.  Existing
\texttt{leanok} nodes name compiled project declarations.  Every other node is
conservatively marked \texttt{notready}.
